\section{An explicit sixfold bound over every imaginary quadratic field, after Markman and Perry}
\label{sec:effective-sixfolds}
This section makes the ideal-curve and Fourier-transform construction quantitative over every imaginary quadratic field. The geometric design is that of Markman's detailed preprint \cite[\S\S8--9]{MarkmanSecant}; its qualitative consequence is also stated in his ICM account \cite[Theorem~1.2]{MarkmanICM}. The deformation input is Perry's Theorem~1.2 \cite{Perry}, an unrefereed preprint: a perfect complex with no negative self-Ext, a coherent finite translation--tensor action, injective semiregularity on invariant obstructions, and an algebraic rational $B$-field whose entire twisted character stays Hodge deforms locally as a twisted perfect complex. We use that theorem with its hypotheses as stated. The estimates below do not use either recent CM Hodge theorem.

\begin{theorem}[With Perry's semiregularity theorem]\label{thm:all-field-sixfolds}
Let $d$ be a positive squarefree integer, let $K=\Q(\sqrt{-d})$, and let $(A,L)$ be a split Weil sixfold with compatible polarization, integral $\mathcal O_f$-action and $\Delta=\chi(L)$. With $m_K$ as in Theorem~\ref{thm:principal-quotient},
\[
 \delta_W(A,L)\le8m_K^3 C_{4d} f^3\Delta
             <10^{179}d^{53}f^3\Delta.
\]
The positive integer $C_D$ is defined explicitly in \eqref{eq:effective-degree} below. When $K$ has class number one, the sharper bound is
\[
 \delta_W(A,L)\le8C_{4d}f^3\Delta
                  <10^{178}d^{51}f^3\Delta.
\]
The detecting cycle may be chosen reduced and irreducible.
\end{theorem}

We prove the theorem in four steps. First a genus-three curve supplies quotient ideals whose obstruction trace is injective. Two distinct elliptic directions make their Fourier transform, after duality, a coherent sheaf of projective dimension at most one. An explicit regularity estimate then turns its local deformation into a globally generated vector bundle, whose third Chern class is an effective detector. Finally a fixed isogeny places this local construction in the principal hyperbolic family, and Theorem~\ref{thm:principal-quotient} transports it to every target lattice.

\subsection{The Fermat curve and its quotient ideals}
Let $C\subset\mathbf P^2$ be $x^4+y^4+z^4=0$. It is a smooth nonhyperelliptic genus-three curve. The three sign involutions give maps $f_j:C\to E_i$ of degree two. Each has four fixed points, so its quotient has genus one. One quotient is $u^2=-(t^4+1)$; the automorphism $(t,u)\mapsto(it,u)$ fixes a point and acts by $i$ on its differential, identifying its Jacobian with $E_i$. The Klein-four quotient of $C$ is the conic $X^2+Y^2+Z^2=0$, so the three nontrivial characters exhaust $H^1(C,\Q)$.

Put $X=E_i^3$ and let $H$ be its product principal polarization. The map
\[
 q:X\longrightarrow J(C),\qquad q(z_1,z_2,z_3)=\sum_j f_j^*z_j
\]
satisfies $q^*\Theta=2H$, $\deg q=8$ and $\ker q\subset X[2]$. Indeed the cross norms $f_{j,*}f_k^*$ vanish for distinct Klein characters and the diagonal norms are $[2]$. Thus the polarization matrix is diagonal $[2]$; Riemann--Roch gives the degree, and the norm identity gives the kernel assertion.

Each $f_j^*$ is injective. A nonzero kernel point would be two-torsion and would factor $C\to E_i$ through a connected \'etale double cover, forcing the remaining map to have degree one and contradicting $g(C)=3$. For an Abel image of $C$, one has the stronger identity
\begin{equation}\label{eq:bielliptic-disjointness}
 f_j^*\operatorname{Pic}^0(E_i)\cap(C-C)=\{0\}.
\end{equation}
To prove it, write $p-q=f_j^*M$ and let $\sigma$ be the involution. Then $p+\sigma q\sim q+\sigma p$. On a nonhyperelliptic curve equal degree-two divisor classes have equal effective representatives. If $p\ne q$, equality of these unordered pairs forces both points to be fixed by $\sigma$. The natural linearization on $\mathcal O_C(p-q)$ has fibre sign $-1$ at these two fixed points and $+1$ at the other two: a local parameter and its inverse both change sign. Any other linearization differs by one global character. Neither choice has trivial inertia at all four points, whereas a pullback from the quotient does. This proves \eqref{eq:bielliptic-disjointness}.

Set $D=4d$ and $k=D+1$. For $j=1,2$, choose a point
$\tau_j\in E_i[k]$ of exact order $k$, and let
$u_j:J(C)\to Y_j=J(C)/\langle f_j^*\tau_j\rangle$ be the quotient.
It embeds $C$ as a smooth curve $\bar C_j\subset Y_j$, because its
translates are disjoint by \eqref{eq:bielliptic-disjointness}.
Let $h_j\in\operatorname{NS}(Y_j)_{\Q}$ be the rational polarization
class with $u_j^*h_j=\Theta$. Then
\[
 \operatorname{ch}(I_{\bar C_j})=1-\frac{k}{2}h_j^2+\frac{k}{3}h_j^3.
\]
This follows by pulling back to the disjoint union of $k$ Abel curves, using $[C]=\Theta^2/2$ and $\chi(\mathcal O_C)=-2$.

For completeness, we check the whole obstruction trace. The restriction $H^1(\mathcal O_{Y_j})\to H^1(\mathcal O_C)$ is an isomorphism. In the normal sequence, the dual of $H^1(T_C)\to H^1(T_{Y_j}|_C)$ contains
$\operatorname{Sym}^2H^0(K_C)\to H^0(K_C^2)$. This map is an isomorphism: both dimensions are six and a plane quartic has no quadratic equation. Hence $h^0(N_{C/Y_j})=3$.
The ideal sequence and Serre duality bound $\operatorname{ext}^1(I,I)$ above by $3+3$. Translations give three independent determinant-preserving directions, while tensoring by $\operatorname{Pic}^0$ gives three independent determinant directions. Thus
\[
 \operatorname{ext}^1(I,I)=\operatorname{ext}^2(I,I)=6.
\]

The categorical trace followed by Hochschild--Kostant--Rosenberg identifies evaluation on $I$ with contraction against $\operatorname{ch}(I)$; the precise compatibility is \cite[Proposition~6.2.1, Corollary~6.3.2 and Theorem~6.4.2, author pp.~29--31]{BFHH}. The tangent bundle is trivial, so the Todd correction is one. The degree-two target is
\[
 HH_{-2}(Y_j)=H^2(\mathcal O_{Y_j})\oplus H^3(\Omega^1_{Y_j}),
\]
of dimension six. The first block is reached isomorphically through the rank term. Contraction from $H^1(T_{Y_j})$ with the nonzero multiple $-kh_j^2/2$ surjects onto the second block: contraction with $h_j$ reaches $H^2(\mathcal O_{Y_j})$, and wedging with $h_j$ identifies that space with $H^3(\Omega^1_{Y_j})$. In coordinates $h_j=\sum x_iy_i$, the latter sends $y_1y_2$ to a nonzero multiple of $x_3y_1y_2y_3$, and likewise for the other basis vectors. The trace therefore has rank six and is injective. In degree one the same calculation separates the Picard directions through the rank term and the translation directions through contraction with $h_j^2$. The degree-zero trace is the nonzero Chern character.

External product preserves this injectivity in degree two: the three K\"unneth terms of self-Ext map to the distinct Hochschild bidegrees $(2,0)$, $(1,1)$ and $(0,2)$, by tensor products of the three injective trace maps just checked. Pullback from a finite quotient identifies invariant self-Ext with the self-Ext downstairs. Trace is natural under that pullback. Thus the pullbacks of these two quotient ideals have a coherent finite descent action and injective invariant obstruction trace.

\subsection{A Fourier transform that is a sheaf after duality}
Put $A=X\times\widehat X$, and write $H_x,H_\alpha$ for the product principal classes on its factors. Write $r_j=u_jq$, $K_j=\ker r_j$, and $F_j=r_j^*I_{\bar C_j}=I_{Z_j}$. Then $|K_j|=8k$; concretely $K_1=\ker q+\langle(\tau_1,0,0)\rangle$ and $K_2=\ker q+\langle(0,\tau_2,0)\rangle$. These summands intersect trivially by injectivity of $f_j^*$. Fourier transform with the normalized Poincar\'e bundle is an equivalence by \cite[Theorem~2.2, p.~156]{Mukai}; its translation, tensor and isogeny identities are in \cite[\S3, especially (3.1) and Proposition~3.11]{Mukai}.
Let $L=\mathcal O_X(2H)$, $\mu(x,y)=(x+y,y)$, and use the inverse
Fourier kernel $\mathcal P^{-1}[3]$ in the second variable. Define
\begin{equation}\label{eq:effective-transform}
 E_{\mathrm{ref}}=
 (\operatorname{id}\times\mathrm{FM}_{\mathcal P^{-1}[3]})
 \mu^*\bigl(([-1]^*F_2\otimes L)\boxtimes(F_1\otimes L)\bigr).
\end{equation}
The order of the factors gives the following incidence description. On $X_y\times A$, put
\[
 a(y,x)=y+x,\quad
 \mathcal L=L_y\otimes a^*L\otimes\mathcal P_{y,\alpha}^{-1},
 \quad D_1=\{y\in Z_1\},\quad D_2=\{y+x\in-Z_2\}.
\]
For $\pi:X_y\times A\to A$ and $T=D_1\cap D_2$, set
\begin{equation}\label{eq:effective-duality}
 G=R\pi_*(\mathcal L I_{D_1}I_{D_2}),
 \qquad E_{\mathrm{ref}}=G[3],\qquad F=G^\vee[-1]
               =E_{\mathrm{ref}}^\vee[2].                 
\end{equation}
Every equality here is a derived-functor identification.

Here is the additional geometric fact that makes $F$ a sheaf.
The curve $C'=q^{-1}C$ is connected: the Abel map induces the
surjection $\pi_1(C)\to\pi_1(J(C))$, hence has transitive monodromy
on the fibre of the connected etale cover $q$. It is an etale
degree-eight cover of $C$, and therefore has genus $17$. Projection
of $C'$ along the first elliptic factor is injective. Indeed, two
points in one such fibre differ by a first-factor point; their
images in $C$ differ by an element of $f_1^*\operatorname{Pic}^0(E_i)$, so \eqref{eq:bielliptic-disjointness} makes those
images equal. The difference then lies in
$\ker q\cap(E_i\times0\times0)=0$, since $f_1^*$ is injective.
The same argument applies to the second factor.

Write $\eta_j$ for the pullback of the point class from the $j$th
elliptic factor of $X$, so $H=\eta_1+\eta_2+\eta_3$.
The two projected images are irreducible divisors $B_1,B_2$ on
the corresponding elliptic surfaces. Their projection maps have
degree one, so the inverse-image divisors $S_1,S_2$ on $X$
have classes
\begin{equation}\label{eq:effective-divisor-classes}
 [S_1]=4(\eta_2+\eta_3),\qquad
 [S_2]=4(\eta_1+\eta_3).                                  
\end{equation}
This follows by pushing forward
$[C']=q^*(\Theta^2/2)=2H^2$ along each elliptic projection.
There is no unaccounted projection degree. One may also check
smoothness: the projected divisor has self-intersection $32$ and
arithmetic genus $17$, equal to the genus of its degree-one
normalization $C'$, so the normalization has no singularity defect.

The $k$ translates defining $Z_1$ lie on $S_1$. On every generic
first-factor fibre of $S_1$ they supply $k$ distinct points.
The $k$ reflected translates in $-Z_2-x$ have the analogous
property on the second-factor fibres of $-S_2-x$. A section of
$\mathcal O_X(4H)\otimes P$ vanishing on both unions, where
$P\in\operatorname{Pic}^0(X)$ is arbitrary, must vanish on each
of these entire elliptic fibres: its restriction has degree four,
whereas $k=4d+1\ge5$. Hence its divisor contains both $S_1$ and
$-S_2-x$. They are distinct for every $x$, because their classes
in \eqref{eq:effective-divisor-classes} differ. The residual effective divisor would have class
\[
 4H-[S_1]-[S_2]=-4\eta_3.
\]
Its intersection with the ample class $H^2$ would be $-8$, which
is impossible for an effective divisor. Thus that space of
sections is zero, for every $x$ and every flat factor $P$.

\paragraph{An explicit derived-to-coherent calculation.}
The coordinates $(y,z=y+x)$ identify the two curve-union ideals
with pullbacks from independent factors. They are Tor independent,
and their product equals their intersection. Therefore
\begin{equation}\label{eq:effective-ideal-sequence}
 0\to\mathcal L I_{D_1}I_{D_2}\to\mathcal L
 \to\mathcal L|_{D_1}\oplus\mathcal L|_{D_2}
 \to\mathcal L|_T\to0                                    
\end{equation}
is exact on the total space.
On a fibre of $\pi$, $\mathcal L$ has class $4H$, with
$h^0=64$ and no higher cohomology. Its restriction to a component
$C'$ has degree $48$. Since $g(C')=17$, its first cohomology is zero
and its space of sections has dimension $32$. Consequently
\[
 W=\pi_*\mathcal L,\qquad
 U=\pi_*(\mathcal L|_{D_1}\oplus\mathcal L|_{D_2})
\]
are vector bundles of ranks $64$ and $64k$, respectively.

The map $f:T\to A$ is finite flat over each of its smooth divisor
images. To see the precise degree, ignore translations and write
one component as $C'\times C'\times\widehat X$, with its map to
$X$ given, up to sign, by addition. The diagonal action
$(p',q')\mapsto(p'+h,q'-h)$ of $\ker q$ is free. Its quotient is
the base change of $C\times C\to C^{(2)}\simeq\Theta$ by $q$.
The quotient map is etale of degree eight; the latter map is
finite flat of degree two onto the smooth divisor $q^{-1}\Theta$.
Thus the composite has degree $16$. Smoothness of $\Theta$ follows
from the absence of a degree-two pencil on the nonhyperelliptic
genus-three curve. Each summand of
\[
 \mathcal T=f_*(\mathcal L|_T)
\]
is a vector bundle on a smooth divisor, pushed forward to $A$.
It has projective dimension one.

Applying direct image to \eqref{eq:effective-ideal-sequence}, including the preceding higher
cohomology vanishings, represents $G$ by
\begin{equation}\label{eq:effective-three-terms}
 [W\longrightarrow U\longrightarrow\mathcal T]
 \quad\hbox{in degrees }0,1,2.                            
\end{equation}
The section argument proves that $W\to U$ is fibrewise injective.
Hence $Q=U/W$ is a vector bundle of rank $64(k-1)=r$, and
$G=[Q\to\mathcal T]$ in degrees $1,2$. No assertion that fibrewise
$H^2(G)$ vanishes has been made or used.

Finite duality gives
\[
 \mathcal T^D=\mathcal E xt^1_A(\mathcal T,\mathcal O_A)
 =\bigoplus f_{ab,*}
       ((\mathcal L|_{T_{ab}})^{-1}\otimes\omega_{T_{ab}}),
\]
with all other sheaf Ext groups zero. Dualizing the last two-term
complex and shifting by $-1$ gives an ordinary coherent sheaf and
the exact extension
\begin{equation}\label{eq:effective-extension}
 0\longrightarrow Q^*\longrightarrow F
   \longrightarrow\mathcal T^D\longrightarrow0.          
\end{equation}
Both the local extension and its finite-duality description prove
$\operatorname{pd}F\le1$; local freeness outside a prescribed
codimension is unnecessary.

\paragraph{The finite action and its trace.}
In \eqref{eq:effective-transform}, the first translation parameter belongs to $K_2$ and
the second to $K_1$. With the convention $t_x^*$ for translation,
the transported action on $E_{\mathrm{ref}}$ has point map
\begin{equation}\label{eq:effective-action}
 (x,y)\longmapsto
 \bigl((x-y,2(x+y)),\,(-2x,y)\bigr)
 \quad\hbox{in }A\times\widehat A.                       
\end{equation}
Tensoring by $L$ sends a translation parameter $x$ to translation $x$ with tensor parameter $-2x$. Pullback by $\mu$ sends the two translation coordinates to $(x-y,y)$ and the tensor coordinates to $(-2x,-2(x+y))$. The inverse Fourier kernel exchanges translation $y$ with tensor $P_y$, and tensor $P_b$ with translation $-b$. This proves the displayed map. Its last coordinate recovers $y$ and its first then recovers $x$, so it is injective, with group order $64k^2$. Its coherence is transported from finite descent. Duality replaces a translation/tensor
point $(a,b)$ by $(a,-b)$; invert the dual linearizations to obtain
the coherent action on $F$. The trace is preserved up to invertible
signs under this operation. Explicitly, derived transpose identifies
the degree-two self-Ext spaces, the Atiyah class of the dual is
the negative transpose, and graded cyclicity of the trace gives
the sign $(-1)^p$ on
$\operatorname{tr}(\xi\operatorname{At}^p/p!)$.
The reversal signs for the cohomological and differential-form
factors cancel. Thus the complete trace kernels identify, including
the invariant subspaces. The even shift $[2]$ contributes sign $+1$.
This proves the claimed strong equivariant semiregularity of $F$.
Negative self-Ext groups vanish because $F$ is a sheaf; they also
vanish through the displayed equivalence with the original ideals.

\subsection{The entire Chern character and the split Weil action}
Put $s^2=-D$, $\theta_D=4DH_x+H_\alpha$, $B=-P/2$ and $\lambda_D=(D+1)/(2D)$, where $P$ is the first Chern class of the normalized Poincar\'e bundle on $X\times\widehat X$. The factor used before Fourier transform has character
\[
 \operatorname{ch}(F_j\otimes L)=1+2H-2DH^2-\frac{4D}{3}H^3
     =\alpha_+e^{2sH}+\alpha_-e^{-2sH},
 \quad\alpha_\pm=\frac{1\pm1/s}{2}.
\]
Reflection preserves this even character.

Here is the Fourier calculation, including the signs. On one elliptic pair let $a,b$ be the first coordinates, $c,e$ the integrated coordinates, and $u,v$ the dual coordinates; normalize $\int ce=1$. For $s_1,s_2\in\{s,-s\}$ the kernel exponent is
\[
 2s_1(a+c)(b+e)+2s_2ce-(cv-eu).
\]
Its integral is
\[
 e^{2s_1ab}\{2(s_1+s_2)-(2s_1b-v)(u-2s_1a)\}.
\]
This follows by collecting the coefficient of $ce$ in the exponential. Multiplying by $e^{-(av-bu)/2}$ and expanding the remaining nilpotent terms gives
\[
 \begin{cases}
 4s_1\exp(s_1(4Dab+uv)/(4D)),&s_1=s_2,\\
 \Omega_{s_1}:=4Dab+2s_1(av-bu)-uv,&s_1=-s_2.
 \end{cases}
\]
Take the product over the three elliptic pairs. The shift $[3]$ in the inverse Fourier transform contributes a minus sign. The equal-sign terms give the polynomial
\begin{equation}\label{eq:effective-polynomial}
 \begin{split}
 P_D(z)={}&64D-8(D-1)z-2z^2+\frac{D-1}{12D}z^3\\
 &+\frac{z^4}{96D}-\frac{D-1}{3840D^2}z^5
       -\frac{z^6}{46080D^2}.
 \end{split}
\end{equation}
The opposite-sign terms give $\lambda_Dt_D$, where
\[
 t_D=-\frac12\left(\prod_{j=1}^3\Omega_{s,j}
                      +\prod_{j=1}^3\Omega_{-s,j}\right).
\]
Duality changes the sign of the codimension-$j$ character by $(-1)^j$; the shift $[2]$ has sign $+1$. Consequently the sheaf $F$ satisfies the full identity
\begin{equation}\label{eq:effective-character}
 e^{-B}\operatorname{ch}(F)=-\lambda_Dt_D+P_D(-\theta_D).
\end{equation}

The rational action of $s$ on each elliptic pair is
$J_D=\left(\begin{smallmatrix}0&1/2\\-2D&0\end{smallmatrix}\right)$.
It squares to $-D$, is skew-adjoint for $\theta_D$, and has signature $(3,3)$. Since
$\Omega_s=-(u-2sa)(v-2sb)$, its conjugate products are the two Weil determinant lines. The class $t_D$ is primitive integral: the odd powers of $s$ cancel and the coefficient of $\prod_j u_jv_j$ is one. Direct multiplication gives
\[
 \theta_Dt_D=0,\qquad \int t_D^2=-2048D^3,
 \qquad \chi(\theta_D)=64D^3.
\]
The Hermitian space is split. For example, the two real cycles in the first elliptic factor form a $K$-basis of each pair, with zero diagonal Hermitian entries and nonzero conjugate off-diagonal entries; rescaling one vector gives $U$.

Set $r=64D$ and $\mathcal H=3(H_x+H_\alpha)$. We prove the following stronger form of $560$-regularity: for every $m\ge1$ and every line bundle $M$ such that $c_1(M)-560m^2\mathcal H$ is nef,
\begin{equation}\label{eq:effective-vanishing}
 H^j(A,[m]^*F\otimes M)=0\qquad(j>0).
\end{equation}
\subsection{A regularity estimate valid after every multiplication isogeny}
All the following vanishings allow arbitrary flat factors.
Relative duality and Serre duality show, for every integer $q\ge1$,
\begin{equation}\label{eq:effective-duality-vanishing}
 H^j(A,W^*(q\mathcal H))=H^j(A,U^*(q\mathcal H))=0
 \qquad(j>0).                                             
\end{equation}
Here is the numerical positivity test, so that \eqref{eq:effective-duality-vanishing} does not
depend on a qualitative regularity assertion. Project the duality
calculation onto the fixed $y$ coordinate, or onto the fixed curve
coordinate for a summand of $U$. The line bundle on the remaining
six-dimensional fibre has matrix
\[
 \begin{pmatrix}2-3q&0\\0&-3q\end{pmatrix}
 \quad\hbox{or}\quad
 \begin{pmatrix}2-3q&1\\1&-3q\end{pmatrix}.
\]
Both are negative definite: the largest eigenvalue before
subtracting $3qI$ is at most $1+\sqrt2<3$. Only the top direct
image of this antiample line bundle can be nonzero, which gives
the required cohomology vanishing below degree six on the dual
side. This proves \eqref{eq:effective-duality-vanishing}. In particular $W^*$ and $U^*$ are
$7$-regular.

Take the evaluation sequence
\[
 0\to K_U\to\mathcal O_A(-7\mathcal H)^N\to U^*\to0.
\]
Its kernel is locally free and $14$-regular. For degrees at least
two this follows directly from \eqref{eq:effective-duality-vanishing} and positivity of
$(7-j)\mathcal H$. For degree one, the multiplication map from
the degree-seven sections, followed by six hyperplanes in the
embedding defined by $\mathcal H$, is surjective by regularity.
Compose the evaluation map with $U^*\to W^*$ and call the kernel
$J$. Then
\begin{equation}\label{eq:effective-kernel-sequence}
 0\to K_U\to J\to Q^*\to0.                              
\end{equation}
The Buchsbaum--Rim resolution of $J$ has terms
\begin{equation}\label{eq:effective-br}
 C_j=\bigwedge^{64+j}\mathcal O_A(-7\mathcal H)^N
          \otimes\operatorname{Sym}^{j-1}W\otimes\det W,
 \qquad j\ge1.                                           
\end{equation}
The first term surjects onto $J$. Local exactness follows by
splitting the surjection as $K\oplus V\to V$ and decomposing
according to the number of factors from $K$. The resulting
strands are the polynomial de Rham complex in positive degree,
contracted by the Euler vector field in characteristic zero;
the degree-one strand maps isomorphically to $K$.

For an explicit bound on \eqref{eq:effective-br}, the theorem of the square gives
\[
 W=\mathcal O_A(2H_x)\otimes b^*\mathrm{FM}(\mathcal O_X(4H)),
 \quad b=2x-\alpha,
 \quad [4]^*W=Q_0^{\oplus64},
\]
with
\begin{equation}\label{eq:effective-four-pullback}
 c_1(Q_0)=16H_x-4H_\alpha+8P,\qquad
 \operatorname{Mat}(Q_0)=\begin{pmatrix}16&8\\8&-4\end{pmatrix}.
\end{equation}
Indeed the pullback of the Fourier bundle by the polarization
isogeny of $\mathcal O(4H)$ is
$H^0(\mathcal O(4H))\otimes\mathcal O(-4H)$; the displayed
formula is its pullback by $b$ with the first line-bundle factor.
The eigenvalues of the matrix in \eqref{eq:effective-four-pullback} are $6\pm2\sqrt{41}$,
whose minimum is greater than $-7$, while
$[4]^*\mathcal H$ has matrix $48I$. The symmetric and determinant
factor in \eqref{eq:effective-br} pulls back to copies of $Q_0^{j+63}$.
Kodaira vanishing and the trace splitting of the isogeny therefore
give
\begin{equation}\label{eq:effective-term-regularity}
 \operatorname{reg}_{\mathcal H}
   (\operatorname{Sym}^{j-1}W\otimes\det W)\le j+69,
 \qquad \operatorname{reg}C_j\le517+8j.                  
\end{equation}
Only six terms are needed to bound $J$ on a sixfold: successive
cohomology sequences leave a final kernel in cohomological
degree seven. Thus
\begin{equation}\label{eq:effective-reg560}
 \operatorname{reg}J\le
 \max_{1\le j\le6}(\operatorname{reg}C_j-j+1)
 \le560,\qquad\operatorname{reg}Q^*\le560.               
\end{equation}

For the torsion dual, use $\rho:C'\times C'\to C\times C$.
On $C\times C$ write $F_1,F_2$ for the point fibres,
$S=F_1+F_2$, and $\Delta_C$ for the diagonal. Addition pulls
$\Theta$ back to $A_C=4S-\Delta_C$. The divisor
$R_C=3S-\Delta_C$ is nef. A degree-three map $C\to\mathbb P^1$
from plane projection supplies its effective residual fibre
correspondence. If its monodromy is $S_3$, that residual divisor
is irreducible and has square $2$, hence is nef. If its monodromy
is cyclic of order three, it is the sum of the two automorphism
graphs: each has square $-4$, and they intersect transversely
in the five fixed points given by Riemann--Hurwitz. Their sum
therefore has intersection $1$ with each component and nonnegative
intersection with every other curve, again proving nefness.

Put $A'=\operatorname{add}_X^*H=\rho^*A_C/2$ and
$H'_i=\rho^*(3F_i/2)$. Up to flat factors on each translated
component,
$c_1(\mathcal L_T)=2H'_1+2H'_2-P_{1,\alpha}$ and
$f^*\mathcal H=3A'+3H_\alpha$. For $q\ge4$ decompose
\begin{equation}\label{eq:effective-torsion-positive}
 \begin{split}
 qf^*\mathcal H-c_1(\mathcal L_T)
 ={}&[3qA'-3H'_1-2H'_2]+(3q-1)H_\alpha\\
    &+[H'_1+H_\alpha+P_{1,\alpha}].
 \end{split}                                               
\end{equation}
The last bracket is a nef addition pullback. The first equals
\[
 \rho^*\left[\frac{3q}{2}R_C+
       \left(\frac{3q}{2}-\frac92\right)F_1+
       \left(\frac{3q}{2}-3\right)F_2\right]
\]
and is ample. Consequently \eqref{eq:effective-torsion-positive} is ample on the smooth
fivefold $C'\times C'\times\widehat X$. Kodaira vanishing and
finite duality give $9$-regularity of $\mathcal T^D$.
Together with \eqref{eq:effective-extension} and \eqref{eq:effective-reg560}, this proves the stated
$560$-regularity of $F$.

For the stronger assertion \eqref{eq:effective-vanishing}, pull the exact constructions
back by $[m]$. The matrices proving \eqref{eq:effective-duality-vanishing} are multiplied by
$m^2$, so the assumed nef difference makes the duality
vanishings hold after tensoring by $M$. The pulled evaluation
sequence has ample middle terms and gives vanishing of
$H^j([m]^*K_U\otimes M)$ for $j\ge2$.
After one further pullback by $[4]$, the first six terms of
\eqref{eq:effective-br} are sums of line bundles whose matrices, after removal
of a nef class, are
\[
 m^2\{48(112-7j)I+(j+63)\operatorname{Mat}(Q_0)\},
 \qquad1\le j\le6.
\]
Their least eigenvalue is at least
$m^2(3360-7\cdot69)=2877m^2>0$. The same six-term truncation
and \eqref{eq:effective-kernel-sequence} give the desired vanishing for $[m]^*Q^*\otimes M$.
For the torsion term, finite flat base change gives a smooth
etale cover $T'$ of $T$; the required line bundle is the sum
of the pullback of \eqref{eq:effective-torsion-positive} with $q=4$ and a nef class.
Kodaira vanishing on $T'$ proves the last vanishing.
The extension \eqref{eq:effective-extension} proves \eqref{eq:effective-vanishing}.

\paragraph{Normalization with an explicit regularity integer.}
Keep $D=4d$, $r=64D$ and the sheaf $F$ just constructed. Set
\begin{equation}\label{eq:effective-normalization}
\begin{gathered}
 h=3\theta_D,\quad \mathcal D=\det F,\quad
 a=\frac{D-1}{8D},\\
 R=560r^2+8Dr+6=2294272D^2+6,\qquad
 B_0=[r]^*F\otimes\mathcal D^{-r}.
\end{gathered}                                             
\end{equation}
The notation $B_0$ here denotes the ordinary normalized sheaf;
the rational B-field used below is written $\mathcal B_0$.
The class $h-\mathcal H$ is nef. Formula \eqref{eq:effective-character} gives
$c_1(F)=-32DP+8(D-1)\theta_D$, and the class
$8Dh-c_1(F)$ has matrix
\[
 \begin{pmatrix}32D(2D+1)&32D\\32D&8(2D+1)\end{pmatrix},
 \qquad\det=256D(4D^2+1)>0.
\]
It is ample. For $1\le i\le6$, the line bundle
$\mathcal D^{-r}\otimes h^{R-i}$ therefore satisfies
\[
 c_1(\mathcal D^{-r}\otimes h^{R-i})-560r^2\mathcal H
 =560r^2(h-\mathcal H)+r(8Dh-c_1(F))+(6-i)h,
\]
which is ample. Assertion \eqref{eq:effective-vanishing} proves that $B_0$ is
$R$-regular. It still has projective dimension at most one.
Its complete character is
\begin{equation}\label{eq:effective-normal-character}
 \operatorname{ch}(B_0)=e^{-r^2a\theta_D}
       \{-\lambda_Dr^6t_D+P_D(-r^2\theta_D)\}.            
\end{equation}
In particular $c_1(B_0)=0$, and, writing its first two relevant
scalar coefficients relative to $h$,
\begin{equation}\label{eq:effective-a2a3}
 \begin{aligned}
 \operatorname{ch}_2(B_0)&=a_2h^2,
 &a_2&=-\frac{r^4(D+1)^2}{18D},\\
 \operatorname{ch}_3(B_0)&=-\lambda_Dr^6t_D+a_3h^3,
 &a_3&=\frac{r^6(D-1)(D+1)^2}{648D^2}.
 \end{aligned}                                             
\end{equation}

\paragraph{Deformation and the ordinary sheaf family.}
Apply Perry's Theorem~1.2 \cite{Perry} to the sheaf $F$, the
coherent finite action proved above, and the algebraic
rational B-field $\mathcal B_0=-c_1(F)/r$. The whole twisted
character is the normalized expression in \eqref{eq:effective-normal-character} before the
$[r]$ scaling: it is a polynomial in the fixed polarization
plus a nonzero Weil class. It is flat and remains of Hodge type
through the balanced Weil family. The determinant monodromy is an integral unit of $K$ of norm one and hence a root of unity; a finite cover makes the Weil class flat.

For precision about untwisting, Perry's \S6.2, proof of Theorem~1.2, pp.~26--28, begins with the $\mu_n$-gerbe representing
the B-field denominator. Here take $n=r$. Its subsequent pushout
to $\mu_m$ does not change the associated multiplicative-group gerbe. Pullback along the map from the original $\mu_r$-gerbe to this pushout is an equivalence on the weight-one categories: locally both are modules on which the band acts through its standard character, and the descent cocycle is unchanged. We therefore pull the deformed weight-one object back to the original finite-band gerbe before untwisting. On a
contractible analytic base the original class remains in the
local system $H^2(A,\mu_r)$. Multiplication by $r$ acts as
$r^2$ on this group and therefore kills it. Pullback by $[r]$
and an analytic Cech trivialization give an ordinary perfect
family. Here is its exact seed comparison, including the sign.
The seed gerbe is the $r$th-root gerbe of $\mathcal D^{-1}$;
its tautological weight-one line $\mathcal T_0$ satisfies
$\mathcal T_0^r=\mathcal D^{-1}$, and the seed object is
$F\otimes\mathcal T_0$. Write its class without an exponential
convention as the Kummer boundary
$\vartheta_0=\delta(\mathcal D^{-1})\in H^2(A,\mu_r)$.
Its rational B-field is
$c_1(\mathcal T_0)=c_1(\mathcal T_0^r)/r
=-c_1(\mathcal D)/r=\mathcal B_0$.
After pullback by $[r]$ and neutralization, the line becomes
an ordinary line $N_0$ with
$N_0^r=[r]^*\mathcal D^{-1}$. Hence
$c_1(N_0)=-r c_1(\mathcal D)$ and
$N_0=\mathcal D^{-r}\otimes P_0$ for some
$P_0\in\operatorname{Pic}^0(A)$. The ordinary seed is therefore
exactly $B_0\otimes P_0$, as asserted.
For these conventions see \cite[Definition~4.5, Lemma~4.6, Remark~4.13, Definition~4.16 and Theorem~4.17, author pp.~15--18]{Hotchkiss}. The displayed root and tensor identities give the required
comparison directly. All preceding estimates allow the flat
factor $P_0$.

The locus where the derived fibres of a relative perfect
complex have cohomology only in degree zero is open. At the
seed this property holds by \eqref{eq:effective-extension}; after shrinking the base
the ordinary family is therefore a coherent sheaf flat over
the base. This can equally be seen by the relative Tor-amplitude
criterion applied to a local finite vector-bundle resolution.
Properness permits a common neighbourhood for the whole fibre.
Analytic upper semicontinuity preserves the six regularity
vanishings of $B_0$ on a smaller neighbourhood. No stability
of all sheaves with this character is used.

\paragraph{A vector-bundle kernel supplies the effective cycle.}
At the seed and throughout that neighbourhood, evaluation gives
\begin{equation}\label{eq:effective-evaluation}
 0\longrightarrow V\longrightarrow
    \mathcal O_A(-Rh)^N\longrightarrow B_0\longrightarrow0,
 \qquad N=h^0(B_0(Rh))=\chi(B_0(Rh)).                     
\end{equation}
At the seed $V$ is locally free, because $\operatorname{pd}B_0\le1$.
This follows locally from the exact Ext sequence, or by comparing
the evaluation presentation with a projective resolution of
length one. In the family the direct image of $B_0(Rh)$ is a
rank-$N$ vector bundle, by the regularity vanishings and base
change; evaluation is surjective. Its kernel is flat, and its
locally free locus is open. The closed complement misses the
proper seed fibre, so it misses every fibre after one more
shrinking. Thus \eqref{eq:effective-evaluation} is a relative vector-bundle
presentation on the neighbourhood.

The kernel $V$ is $(R+7)$-regular. For cohomological degrees
at least two, use the sequence \eqref{eq:effective-evaluation}, the regularity of $B_0$,
and the positive line bundles $(7-i)h$ for $2\le i\le6$.
For degree one, $R$-regularity gives the surjection
\[
 H^0(B_0(Rh))\otimes\operatorname{Sym}^6H^0(\mathbb P^M,\mathcal O(1))
 \longrightarrow H^0(B_0((R+6)h)),
\]
which factors through the degree-six sections in \eqref{eq:effective-evaluation}.
It gives the remaining $H^1$ vanishing. Therefore
$V'=V((R+7)h)$ is a globally generated vector bundle.

Let $q_0=R+7$ and define the explicit rational numbers
\begin{equation}\label{eq:effective-uvw}
 \begin{split}
 u&=7N-rq_0,\\
 v&=\frac{49N}{2}-\frac{rq_0^2}{2}-a_2,\\
 w&=\frac{343N}{6}-\frac{rq_0^3}{6}-q_0a_2-a_3.
 \end{split}                                               
\end{equation}
The identity in $K_0(A)$ is
\[
 [V']=N[\mathcal O_A(7h)]-[B_0(q_0h)].
\]
Consequently $c_1(V')=uh$, $\operatorname{ch}_2(V')=vh^2$,
and $\operatorname{ch}_3(V')=wh^3+\lambda_Dr^6t_D$.
The exact Newton identity
$c_3=2\operatorname{ch}_3-c_1\operatorname{ch}_2+c_1^3/6$
therefore gives
\begin{equation}\label{eq:effective-c3}
 c_3(V')=2\lambda_Dr^6t_D+
                (2w-uv+u^3/6)h^3.                       
\end{equation}
Its Weil projection is nonzero. In particular the rank of $V'$
is at least three, since lower-rank bundles have $c_3=0$.

General sections of a globally generated vector bundle give
an effective maximal-rank-drop locus of expected codimension
three, or the empty locus, representing $c_3(V')$. This is the
elementary Thom--Porteous construction: take
$\operatorname{rank}(V')-2$ sections, use the universal surjective
evaluation to get transversality to the rank strata, and then
take a general choice. Equation \eqref{eq:effective-c3} rules out the empty case.
It is an effective cycle with integral coefficients with nonzero Weil projection,
so at least one reduced irreducible component detects the Weil
plane. No condition on the codimension of the nonfree locus of
$B_0$ is required.

Write $P_D(z)=\sum_{j=0}^6p_jz^j$, with coefficients given in
\eqref{eq:effective-polynomial}. Since the Todd class of an abelian variety is $1$ and
$\theta_Dt_D=0$, the integer $N$ is explicitly
\[
 N=(6!\,64D^3)\sum_{j=0}^6
       p_j(-r^2)^j\frac{(3R-r^2a)^{6-j}}{(6-j)!}.
\]
The exact degree of the effective cycle above is the nonnegative
integer
\begin{equation}\label{eq:effective-degree}
 C_D=(3^6\,6!\,64D^3)(2w-uv+u^3/6).                     
\end{equation}
The formula defines every quantity by rational operations on
the specified positive integer $D$. Integrality and nonnegativity
also follow from the constructed vector bundle and effective
cycle. These same Chern classes and degrees persist on the
analytic neighbourhood. GAGA makes each fibre and cycle algebraic.

\subsection{The degree and its transport to every field and lattice}
The formulas above specify $C_D$ by rational arithmetic in $D$. We first give a convenient upper bound. Put $s_0=3R-r^2a$. From \eqref{eq:effective-polynomial}, $j!|p_j|\le64D$ for $0\le j\le6$, while $0<s_0$ and $s_0+r^2<7\cdot10^6D^2$. The Euler-characteristic formula therefore yields
\[
 0<N\le4096D^4(s_0+r^2)^6<10^{45}D^{16}.
\]
Also $q_0<3\cdot10^6D^2$, $|a_2|<4\cdot10^6D^5$ and $|a_3|<5\cdot10^8D^7$. Substitution in \eqref{eq:effective-uvw} gives
\[
 |u|\le8\cdot10^{45}D^{16},\quad
 |v|\le26\cdot10^{45}D^{16},\quad
 |w|\le59\cdot10^{45}D^{16}.
\]
As $10^{45}D^{16}\ge1$, the absolute value of $2w-uv+u^3/6$ is at most $512\cdot10^{135}D^{48}$. Hence
\[
 0<C_D<10^{146}D^{51},\qquad C_{4d}<10^{177}d^{51}.
\]
Positivity follows also from the nonzero effective detecting cycle, and integrality follows from its degree.

We next give a fixed map from the source to a principal hyperbolic lattice. Put $s=\sqrt{-d}$ and let $e_d=1$ for $d\equiv1,2\pmod4$, and $e_d=2$ for $d\equiv3\pmod4$. Thus $\tau=s$ or $(1+s)/2$ generates $\mathcal O_K$, and $\tau-\bar\tau=2s/e_d$. In one elliptic pair write the real homology basis as $u=e_x$, $v=i e_x$, $w=e_y$, $z=i e_y$. The $s$-action is
\[
 su=-4dw,\quad sv=-4dz,\quad sw=u/4,\quad sz=v/4.
\]
The source polarization $h=3(16dH_x+H_y)$ has $E_h(u,v)=48d$ and $E_h(w,z)=3$. In the $K$-basis $u,v$ its normalized Hermitian matrix, linear in the first variable, is
\[
 \begin{pmatrix}0&48ds/e_d\\-48ds/e_d&0\end{pmatrix}.
\]
In the principal hyperbolic lattice $\mathcal O_K e\oplus\mathcal O_K f$, define
\[
 T(u)=4s e,\quad T(v)=12d f,\quad
 T(w)=e,\quad T(z)=-3s f.
\]
These values are integral and compatible with the $K$-action, and $H_U(Tu,Tv)=48ds$. Thus $T^*E_U=e_dE_h$. Three copies give an isogeny $\varphi:A_0\to P_U$ with $\varphi^*\Theta_U=e_dh$. In particular the pushforward of a three-cycle costs exactly $e_d^3$ in total cycle degree. Its Weil projection stays nonzero. This calculation uses no class-number assumption.

The fixed rational similitude identifies the period domains. After choosing a common finite level it sends the analytic neighbourhood where the sheaf construction works onto an analytic open subset of the principal $U^3$ family. Bounded-detector closedness therefore gives a detector of degree at most $e_d^3C_{4d}$ on every member of that family. Theorem~\ref{thm:principal-quotient}, and $e_dt_K=2$, now give
\[
 \delta_W(A,L)\le(t_Kfm_K)^3\Delta\,e_d^3C_{4d}
                 =8m_K^3C_{4d}f^3\Delta.
\]
Since $m_K\le4\sqrt d/\pi$, one has $8m_K^3<17d^{3/2}$. The asserted coarse bound follows from $C_{4d}<10^{177}d^{51}$ and $d^{105/2}\le d^{53}$. For class number one, $m_K=1$. A detecting irreducible component never has greater degree than the effective cycle containing it. This proves Theorem~\ref{thm:all-field-sixfolds}.

The exterior-algebra and integer-arithmetic calculations are independently reproducible with the programs \texttt{sixfold\_character\_check} and \texttt{sixfold\_degree\_check} in the \href{\workbenchurl}{workbench}. The proofs above work for every positive squarefree $d$. The programs check the formulas and signs, and do not replace the geometric deformation argument.
